(1) 该场满足 \(\nabla\cdot\mathbf B=G-G=0\)，模长 \(|\mathbf B|=\sqrt{B_0^2+G^2(x^2+y^2)}\)。对弱场寻求态，绝热势 \(U(r)=\mu_{\rm eff}|\mathbf B|\)，其最小值在 \(r=0\)，与强场寻求态的符号相反。
(2) 对 \(Gr\ll B_0\)，\(U(r)=\mu_{\rm eff}B_0+\mu_{\rm eff}G^2r^2/(2B_0)+O(r^4)\)。与 \(mr^2\omega_r^2/2\) 比较，\(\omega_r=G\sqrt{\mu_{\rm eff}/(mB_0)}\)；两个横向自由度简并，轴向无束缚。
(3) 阱口相对最低点的深度 \(\Delta U=\mu_{\rm eff}[\sqrt{B_0^2+G^2a^2}-B_0]\)。仅考虑径向自由度且忽略碰撞时，束缚要求径向总能量低于 \(\Delta U\)；该深度不是原子全三维俘获深度。
(4) 自旋Larmor进动角频率应远大于磁场方向沿原子轨道的变化率，典型条件 \(\omega_L(B_0)\gg\omega_r\)。若 \(B_0\to0\)，轴心磁场零点导致绝热近似失效和Majorana自旋翻转，故小振动公式不能无条件外推。
